\section*{Capítulo \ref{ch:b1:organomagnesium} --- Los compuestos organomagnesianos}

\begin{solution}{exo:b1:organomagnesium:1}
\ce{CH3CH2Br + Mg -> CH3CH2MgBr}; \ce{C6H5Br + Mg -> C6H5MgBr} (en éter
seco). C--Br: el Br (2.96) es más electronegativo que el C (2.55), y el
carbono es positivo. C--Mg: el Mg (1.31) es menos electronegativo, y el
carbono es negativo.
\end{solution}

\begin{solution}{exo:b1:organomagnesium:2}
Agua: \ce{CH3MgI + H2O -> CH4 + Mg(OH)I}.
Etanol: \ce{CH3MgI + CH3CH2OH -> CH4 + CH3CH2OMgI}.
Ácido etanoico: \ce{CH3MgI + CH3COOH -> CH4 + CH3COOMgI}.
Óxido de deuterio: \ce{CH3MgI + D2O -> CH3D + Mg(OD)I}.
\end{solution}

\begin{solution}{exo:b1:organomagnesium:3}
Metanal: propan-1-ol (primario). Etanal: butan-2-ol (secundario).
Propanona: 2-metilbutan-2-ol (terciario). Benzaldehído:
1-fenilpropan-1-ol (secundario). Dióxido de carbono: ácido propanoico.
\end{solution}

\begin{solution}{exo:b1:organomagnesium:4}
Su grupo OH es lo bastante ácido para destruir el reactivo en cuanto se
forma: la molécula se protonaría a sí misma (o a sus vecinas), dando el
alcóxido de magnesio del butan-1-ol, \ce{CH3(CH2)3OMgBr}, que no es un
reactivo utilizable. Antes hay que proteger el OH
(\cref{ch:b1:acetals-protection} muestra la idea para un aldehído).
\end{solution}

\begin{solution}{exo:b1:organomagnesium:5}
El par C--Mg ataca el carbono carbonílico del etanal, y el par $\pi$ del
C=O pasa al O: \ce{CH3CH(CH3)O^-} con \ce{MgBr+}; después, \ce{H3O+} cede
un protón al oxígeno: propan-2-ol. El propan-2-ol lleva dos grupos metilo
en el carbono carbinólico: es \omterm{def:b1:stereochemistry-in-depth:stereocentre}{aquiral}. (Con un R distinto, por ejemplo
el bromuro de etilmagnesio con el etanal, el producto, el butan-2-ol, es
\omterm{def:b1:stereochemistry-in-depth:stereocentre}{quiral} pero racémico: el carbonilo plano es atacado por sus dos caras por
igual.)
\end{solution}

\begin{solution}{exo:b1:organomagnesium:6}
El carbono carbinólico lleva un metilo y dos etilos: bromuro de
etilmagnesio con butan-2-ona, o bromuro de metilmagnesio con
pentan-3-ona.
\end{solution}

\begin{solution}{exo:b1:organomagnesium:7}
Se prepara \ce{(CH3)2CHMgBr} a partir del 2-bromopropano y magnesio en
éter seco, se vierte sobre \ce{CO2} sólido y se acidifica:
\ce{(CH3)2CHCOOH}. $n = 12.3/123.0 = 0.100$ mol;
$0.100 \times 0.70 \times 88.0 = \qty{6.2}{g}$.
\end{solution}

\begin{solution}{exo:b1:organomagnesium:8}
$0.18/18.0 = \qty{0.010}{mol}$ de agua destruyen \qty{0.010}{mol} de
reactivo: un \qty{20}{\percent}; se forma benceno,
\ce{C6H5MgBr + H2O -> C6H6 +
Mg(OH)Br}.
\end{solution}

\begin{solution}{exo:b1:organomagnesium:9}
El reactivo, un \omterm{def:b1:electronic-effects:nucleophile}{nucleófilo}, reacciona con el bromobenceno que aún queda:
en conjunto, \ce{C6H5MgBr + C6H5Br -> C6H5-C6H5 + MgBr2}. Añadir el
bromuro lentamente a un exceso de magnesio mantiene baja su concentración:
encuentra metal limpio antes que el reactivo ya formado.
\end{solution}

\begin{solution}{exo:b1:organomagnesium:10}
1-Feniletan-1-ol: \ce{CH3MgBr} (a partir del bromometano) con
benzaldehído, o \ce{C6H5MgBr} (a partir del bromobenceno) con etanal.
2-Fenilpropan-2-ol: \ce{C6H5MgBr} con propanona, o \ce{CH3MgBr} con
1-feniletan-1-ona.
\end{solution}

\begin{solution}{exo:b1:organomagnesium:11}
La cetona se añade al reactivo para que este, preparado en el matraz
seco, no salga nunca de él, y para que el calor liberado por cada porción
lo absorba el éter en ebullición. El agua sola daría la sal básica
gelatinosa \ce{Mg(OH)Br}, que atrapa el producto; el ácido diluido frío
la disuelve como \ce{Mg^2+}, y el frío limita las reacciones secundarias
del alcohol terciario en medio ácido.
\end{solution}

\begin{solution}{exo:b1:organomagnesium:12}
Cada mol de reactivo da con el agua un mol de sal básica de magnesio,
\ce{RMgBr + H2O -> RH + Mg(OH)Br}, cuyo \ce{OH-} valora el ácido: la
cantidad de reactivo es igual a la de ácido usada, \qty{0.088}{mol}, un
\omterm{def:b1:extent-q-and-k:fractional-extent}{rendimiento} del \qty{88}{\percent}.
\end{solution}

\begin{solution}{pb:b1:organomagnesium:1}
\textbf{1.} \ce{C6H5Br + Mg -> C6H5MgBr}.
\textbf{2.} Mg: $0.535/24.3 = \qty{0.0220}{mol}$; \ce{C6H5Br}: $3.14/156.9 =
\qty{0.0200}{mol}$: el magnesio está en ligero exceso.
\textbf{3.} El agua destruye el reactivo: \ce{C6H5MgBr + H2O -> C6H6 +
Mg(OH)Br}.
\textbf{4.} Impide que el vapor de agua del aire entre por el
refrigerante.
\textbf{5.} El éter cede \omterm{def:b1:lewis-resonance-vsepr:lewis}{pares no enlazantes} al magnesio (\omterm{def:b1:lewis-resonance-vsepr:lewis-acid}{base de Lewis}):
sin él, el reactivo no se forma ni permanece en disolución.
\textbf{6.} \ce{C6H5MgBr + C6H5Br -> C6H5-C6H5 + MgBr2}; bifenilo,
$0.15/154.0 = \qty{9.7e-4}{mol}$, que ha consumido \qty{9.7e-4}{mol} de
bromobenceno y otro tanto de reactivo: en total se han perdido
\qty{1.9e-3}{mol} de bromobenceno.
\textbf{7.} $3.28/182.0 = \qty{0.0180}{mol}$.
\textbf{8.} El par C--Mg ataca el carbono carbonílico y el par $\pi$ pasa
al oxígeno: \ce{(C6H5)3C-O^-} \ce{MgBr+}.
\textbf{9.} Reactivo disponible: como máximo, $0.0200 - 0.0019 = \qty{0.0181}{mol}$;
benzofenona, \qty{0.0180}{mol}: limita la benzofenona, por muy poco.
\textbf{10.} No: el carbono carbinólico lleva tres grupos fenilo
idénticos.
\textbf{11.} El reactivo permanece en su matraz seco, y el calor de cada
porción de cetona lo absorbe el éter en ebullición.
\textbf{12.} Destruiría parte del reactivo (benceno) y dejaría benzofenona
sin reaccionar.
\textbf{13.} \ce{(C6H5)3COMgBr + H3O+ -> (C6H5)3COH + Mg^2+ + Br- + H2O}.
\textbf{14.} El ácido disuelve las sales básicas de magnesio que el agua
sola precipitaría; el frío limita la formación del \omterm{def:b1:electronic-effects:intermediates}{carbocatión} de la
pregunta 18.
\textbf{15.} El trifenilmetanol y el bifenilo, en la capa etérea; las
sales de magnesio, en la capa acuosa.
\textbf{16.} Lavando (o triturando) el sólido con éter de petróleo frío:
el bifenilo se disuelve y el trifenilmetanol queda.
\textbf{17.} Un punto de fusión nítido igual al tabulado (o una sola
mancha en cromatografía en capa fina).
\textbf{18.} \ce{(C6H5)3COH + H+ -> (C6H5)3C+ + H2O}: un \omterm{def:b1:electronic-effects:intermediates}{carbocatión}
terciario cuya carga se reparte por resonancia entre tres anillos de
benceno, muy estable, y coloreado por esa conjugación extendida.
\textbf{19.} $3.75/260.0 = \qty{0.0144}{mol}$.
\textbf{20.} $0.0144/0.0180 = \qty{80}{\percent}$.
\textbf{21.} Pérdidas en los trasvases y en la recristalización; reactivo
destruido por trazas de humedad; reacción incompleta.
\textbf{22.} Ácido benzoico, \ce{C6H5COOH}, tras acidificar.
\textbf{23.} $0.0181 \times 122.0 = \qty{2.21}{g}$.
\textbf{24.} \ce{C6H5MgBr + CO2 -> C6H5COOMgBr} y, después,
\ce{C6H5COOMgBr + H3O+ -> C6H5COOH + Mg^2+ + Br- + H2O}.
\textbf{25.} \textbf{\qty{80}{\percent}}.
\end{solution}
